\documentclass{article}
\pdfpagewidth=8.5in
\pdfpageheight=11in
\usepackage{ijcai27}
\usepackage{times}
\usepackage{url}
\usepackage[hidelinks]{hyperref}
\usepackage[utf8]{inputenc}
\usepackage[small]{caption}
\usepackage{graphicx}
\usepackage{amsmath,amssymb,amsthm}
\usepackage[switch]{lineno}
\linenumbers
\urlstyle{same}
\newtheorem{proposition}{Proposition}
\newcommand{\Hist}{\mathcal{H}}
\newcommand{\Loss}{\mathcal{L}}
\pdfinfo{
/TemplateVersion (IJCAI.2027.0)
}
\title{SAFEPlus: Separating Fraud Occurrence from Delayed Detection}
\author{Anonymous Authors}
\begin{document}
\maketitle
\begin{abstract}
Platforms typically record when an account is suspended rather than when
fraud first occurs. SAFE addresses this mismatch through a survival-based
early-detection objective, but does not explicitly model the delay between
fraud occurrence and platform detection. This research note develops a
candidate SAFEPlus formulation with a latent fraud onset process and a
subsequent detection process. Marginalizing onset time gives an observed-data
likelihood that distinguishes accounts without fraud from accounts with
undetected fraud at the observation boundary. We derive a discrete-time
formulation with exact marginalization and give an elementary example of
non-identifiability from detection times alone. We outline synthetic recovery
experiments and report an exploratory audit and rerun of the legacy SAFE
baseline. SAFEPlus performance, onset recovery, and novelty beyond the
supplied conceptual note remain to be established.
\end{abstract}
\section{Motivation and Scope}
Fraud occurrence and platform intervention are different events. A user may
commit fraud before the platform accumulates sufficient evidence to suspend
the account. Learning suspension time as though it were fraud onset therefore
mixes user behavior with the platform's detection mechanism.

SAFE~\cite{zheng2019safe} combines recurrent sequence modeling with positive
hazard outputs and a monotonically decreasing survival score. Its revised
objective rewards event probability accumulated before an observed suspension,
rather than concentrating event probability at suspension time. This supports
early flagging, but does not explicitly identify when fraud began or how long
the platform took to detect it. Moreover, absence of suspension is not evidence
that fraud has not occurred.

The supplied conceptual note~\cite{smola2022note} proposes separate hazards
for latent fraud occurrence and subsequent detection, together with latent
variable inference. We use that proposal as the starting point of SAFEPlus;
the two-process idea itself is prior conceptual material, not a new
contribution claimed here. Attention, hierarchical representations, and future
distillation described in that note are deferred while the likelihood and its
limitations are examined.

\paragraph{Working research question.}
Under which assumptions can delayed detection records support useful
inference about fraud onset, and when does explicitly modeling detection
delay improve early warning relative to SAFE?

\paragraph{Current scope.}
This draft contains a candidate likelihood, a finite-grid implementation
specification, and an evaluation plan. It reports an exploratory SAFE baseline rerun. A potential
submission contribution must be established through a broader literature
review, explicit identification assumptions, and evidence of reliable
inference or improved decisions. Replacing a GRU with attention alone is not
the present research objective.

\section{A Joint Occurrence--Detection Model}
\subsection{Observations and Assumptions}
For account $i$, let $T_{f,i}$ be its first fraud time and $T_{d,i}$ its first
platform detection time. We observe
\[
 y_i=\min(T_{d,i},C_i),\qquad
 \delta_i=\mathbf{1}\{T_{d,i}\le C_i\},
\]
where $C_i$ is an administrative observation boundary. We also observe a
timestamped covariate history $\Hist_i(t)$ up to $y_i$. Time zero is the start
of the modeled risk period. Initially, all accounts are assumed fraud-free;
prevalent fraud at enrollment requires an additional initial-state mixture.

The initial model assumes that detection follows genuine fraud, with
$T_d\ge T_f$, and that fraud status is absorbing. False suspensions,
repeated fraud episodes, and rehabilitation are outside this first model.
Censoring is assumed non-informative conditional on the modeled information.
These are modeling assumptions to audit, not established properties of the
Twitter or Wiki data.

We first specify the likelihood conditional on an external covariate path;
path conditioning is suppressed in the notation. Neural hazards may use
causal features available immediately before each transition. If covariates
are themselves generated by the latent fraud state, this conditional
construction needs further justification or a joint observation model.
No future observations may be used for an online prediction.

\subsection{Continuous-Time Likelihood}
Let $\lambda_f(t)$ be the hazard of first fraud, and let
$\lambda_d(t\mid u)$ be the detection hazard at $t$ after fraud at $u$.
Dependence on onset $u$ can include elapsed time $t-u$. Define
\begin{align}
 S_f(t)&=\exp\left[-\int_0^t\lambda_f(s)\,ds\right],\\
 f_f(u)&=\lambda_f(u)S_f(u),\\
 S_d(t\mid u)&=\exp\left[-\int_u^t\lambda_d(s\mid u)\,ds\right],\\
 f_d(t\mid u)&=\lambda_d(t\mid u)S_d(t\mid u).
\end{align}
We do not require that all accounts eventually commit fraud. The
occurrence distribution may retain probability mass at infinity.

For a detection observed at $d$, the marginal density is
\begin{equation}
 p_d(d)=\int_0^d f_f(u)f_d(d\mid u)\,du.
 \label{eq:detected}
\end{equation}
For an account undetected at $C$, its probability is
\begin{equation}
 P(T_d>C)=S_f(C)+\int_0^C f_f(u)S_d(C\mid u)\,du.
 \label{eq:censored}
\end{equation}
The terms in Equation~\ref{eq:censored} represent no fraud by $C$, and fraud
by $C$ without detection, respectively. Treating all undetected users as
fraud-free removes the second explanation.
The observed-data negative log-likelihood is
\begin{equation}
 \Loss=-\sum_i\left[\delta_i\log p_d(y_i)
 +(1-\delta_i)\log P(T_{d,i}>y_i)\right].
 \label{eq:nll}
\end{equation}
Exact timestamps contribute densities. If platform status is checked only
periodically, detection is interval-observed and its contribution is instead
$P(l_i<T_{d,i}\le r_i)$. This distinction must be resolved in preprocessing.

\subsection{Discrete-Time Specification}
For an initial implementation, let $k=1,\ldots,K$ index observation bins.
Let $a_k\in(0,1)$ be the probability of first fraud in bin $k$, conditional on
no previous fraud. Given onset in bin $k$, let $b_{j\mid k}\in(0,1)$ be the
conditional detection probability in bin $j\ge k$, given no earlier detection.
The bin convention allows onset followed by detection within the same bin.
This defines a discrete model; it is not an exact continuous-time model
without an additional discretization argument.

The onset mass and conditional detection survival are
\begin{align}
 g_k&=a_k\prod_{r<k}(1-a_r),\\
 Q_{j\mid k}&=\prod_{r=k}^{j}(1-b_{r\mid k}).
\end{align}
An empty product is one. Detection in bin $d$ and non-detection through bin
$c$ have likelihoods
\begin{align}
 p_d(d)&=\sum_{k=1}^{d}g_k Q_{d-1\mid k}b_{d\mid k},
 \label{eq:discrete-detected}\\
 p_{\mathrm{und}}(c)&=\prod_{r=1}^{c}(1-a_r)
                  +\sum_{k=1}^{c}g_k Q_{c\mid k}.
 \label{eq:discrete-censored}
\end{align}
These finite sums permit exact marginalization over onset bins and direct
automatic differentiation, avoiding a sampling approximation in the first
experiment. Computation should use log probabilities and log-sum-exp.
General onset-dependent detection costs $O(K^2)$ per sequence.

\paragraph{A linear-time restricted model.}
If detection depends on current causal features but not onset bin,
$b_{j\mid k}=b_j$, define $U_k$ as the probability of no fraud through bin $k$
and $V_k$ as the probability of fraud but no detection through that bin.
Starting with $U_0=1$ and $V_0=0$,
\begin{align}
 U_k&=U_{k-1}(1-a_k),\\
 V_k&=(V_{k-1}+U_{k-1}a_k)(1-b_k),\\
 D_k&=(V_{k-1}+U_{k-1}a_k)b_k.
\end{align}
Here $D_k$ is the probability mass of first detection in bin $k$.
The censored likelihood is $U_c+V_c$. These recurrences imply
$U_c+V_c+\sum_{k=1}^cD_k=1$, providing a useful implementation check.
A shared causal GRU with separate sigmoid heads is a candidate
parameterization of $a_k$ and $b_k$, not a commitment that this
parameterization resolves identification.

\subsection{Retrospective and Online Inference}
For a detected account, retrospective onset responsibilities are
\begin{equation}
 P(T_f=k\mid T_d=d)
 =\frac{g_k Q_{d-1\mid k}b_{d\mid k}}{p_d(d)}.
\end{equation}
For an undetected account, the posterior includes both an onset-bin
component $g_kQ_{c\mid k}/p_{\mathrm{und}}(c)$ and a no-fraud-yet component
$U_c/p_{\mathrm{und}}(c)$.
In the restricted model, an online score among undetected accounts is
\begin{equation}
 R_c=P(T_f\le c\mid T_d>c,\Hist(c))=\frac{V_c}{U_c+V_c}.
\end{equation}
This filtered probability need not be monotone: continued non-detection can
change the evidence about past fraud. An absorbing alert policy may retain a
previous alert, but that policy is distinct from probability monotonicity.
Retrospective smoothing uses later detection information and must not be
reported as an online prediction made earlier.

\section{What Can Be Identified?}
Accurate prediction of detection does not by itself establish accurate
recovery of fraud onset. The following restricted example already exhibits
ambiguity.

\begin{proposition}
Suppose $T_f\sim\mathrm{Exp}(\alpha)$ and the independent detection delay
$L=T_d-T_f\sim\mathrm{Exp}(\beta)$, with positive unequal rates. The
distribution of $T_d$ alone cannot distinguish $(\alpha,\beta)$ from
$(\beta,\alpha)$.
\end{proposition}
\begin{proof}
The convolution density is
\begin{equation}
 p_d(t)=\frac{\alpha\beta}{\beta-\alpha}
       \left(e^{-\alpha t}-e^{-\beta t}\right).
\end{equation}
Exchanging the rates changes the signs of both the denominator and the
parenthesized difference, leaving the density unchanged. Detection survival
is also unchanged, so administrative right-censoring cannot resolve the
ambiguity. The mean fraud onset, $1/\alpha$, generally changes.
\end{proof}

This is a counterexample, not a general identification theorem. It shows
that detection-time fit alone is insufficient even in a simple submodel.
Flexible neural hazards can introduce further ambiguity. Restricting a
hazard or regularizing parameters selects solutions but does not establish
that the selected onset corresponds to the true process.

Possible additional information includes an audited subset with onset labels
or intervals, externally estimated detection delays, or covariates with
justified effects on only one process. Each route requires explicit
assumptions and sensitivity analysis. We will first compare unconstrained
estimation with known-delay and partially audited settings. A meaningful
outcome may be calibrated onset intervals or sensitivity bounds rather than
point identification.

\section{Evaluation}
\subsection{Exploratory SAFE baseline rerun}
We reran SAFE on the bundled Twitter and Wiki splits using one seed
and TensorFlow 2.12 in compatibility mode. The cumulative-hazard loss
was evaluated with an algebraically equivalent stable expression
after the original expression overflowed on Twitter. Model updates
used training data. Checkpoint and threshold selection maximized mean
validation accuracy over prefixes 1--5. Validation and test inference
received only the first five feature vectors per account. All test
records were scored, including remainder batches.

\paragraph{Split audit.}
We preserved the test splits and removed records from earlier splits
when their full feature sequences matched another split, giving test
priority over validation and validation priority over training.
On Twitter, this removed 64 training and 15 validation records,
leaving 3,806 training, 541 validation, and 1,114 test records.
Wiki had no exact full-sequence matches across splits and retained
1,231 training, 175 validation, and 353 test records. After filtering,
neither dataset had exact full-sequence matches across splits.
Exact matches are a feature-level check. Sparse sequences may match
across different accounts, and the bundled arrays lack account IDs and
source timestamps for verifying account separation and feature availability.

\begin{table}[t]
\centering
\small
\setlength{\tabcolsep}{4pt}
\begin{tabular}{llrrrr}
\hline
Dataset & Run & Prec. & Recall & F1 & Acc. \\
\hline
Twitter & Published & 82.0 & 55.7 & 65.4 & 71.8 \\
        & Audited rerun & 71.8 & 60.0 & 64.8 & 67.9 \\
Wiki    & Published & 71.1 & 88.0 & 78.7 & 76.4 \\
        & Audited rerun & 69.7 & 89.0 & 78.2 & 75.5 \\
\hline
\end{tabular}
\caption{SAFE performance (\%) averaged over prefixes 1--5.
The published study reports ten runs. The audited rerun uses one
seed per dataset.}
\label{tab:safe-audited-rerun}
\end{table}

\paragraph{Results and limitations.}
The published values are from Table~1 of \cite{zheng2019safe}.
In Table~\ref{tab:safe-audited-rerun}, local F1 was within 0.6
percentage points of the published score on both datasets. Twitter
precision remained 10.2 points below the published value.
The rerun changed overlap handling, batching, and the validation
window together. The legacy test results had been inspected before
these changes, so the rerun is exploratory. A confirmatory evaluation
requires a fresh account-disjoint holdout, verifiable feature
timestamps, and a protocol fixed before test evaluation.

\subsection{Planned SAFEPlus evaluation}
\paragraph{Stage 1: likelihood verification.}
Enumerate short sequences to compare Equations~\ref{eq:discrete-detected}
and~\ref{eq:discrete-censored} against the linear-time recurrence. Check
normalization, posterior sums, gradient calculations, variable sequence
lengths, and censoring boundaries. Generate an explicit detection observation
indicator separately from latent fraud status.

\paragraph{Stage 2: synthetic recovery.}
Generate known onset and delay distributions, then hide onset during fitting.
Vary delay length, censoring, sample size, and the fraction of audited onset
labels. Include the rate-swapping example as a deliberate failure case.
Compare no delay information, a known delay process, and partial onset
supervision across multiple initializations. Report held-out detection
log-likelihood separately from onset error, posterior interval coverage,
undetected-fraud calibration, and parameter sensitivity. Strong detection fit
with poor onset recovery is an important diagnostic outcome.

\paragraph{Stage 3: controlled comparison with SAFE.}
Establish the original data and loss conventions before porting the legacy
implementation to a current framework. Compare SAFE, a standard detection-time
survival model, and the proposed joint model under matched causal encoders,
splits, and tuning budgets. Add a supervised-onset oracle only in synthetic
experiments as a reference. Introduce an onset-dependent delay model only
after the restricted model is validated.

\paragraph{Real-data audit and metrics.}
Determine whether timestamps denote observations, edits, suspensions, or
intervals between status checks. Audit label semantics, censoring, padding,
and train/validation/test separation. Select alert thresholds on validation
data only. Report prefix-based precision--recall performance, false alerts,
and lead time relative to observed suspension, with repeated-run uncertainty.
Where chronological splits are possible, prevent future account or platform
information from entering training features. Without independently verified
onset labels, earlier alerts than suspension do not establish accurate fraud
onset recovery or anticipation before fraud occurs.

\section{Limitations and Next Steps}
The initial formulation assumes a single irreversible fraud onset, reliable
detection labels, and an observation mechanism compatible with the stated
censoring model. Platform policies may change, suspicious accounts may be
monitored more intensely, and suspensions may be noisy. The current
conditional treatment of behavioral covariates also needs an explicit audit.
These issues affect interpretation of the hazards and inferred onset times.

Before a submission is defensible, the draft needs a broader related-work
review on delayed feedback, latent event times, multi-state survival models,
and partial supervision. The proposed likelihood must be distinguished from
existing models and from the supplied conceptual note. The origin,
publication status, and appropriate attribution of that note must be
confirmed before submission. Its bibliography entry is provisional.

\paragraph{Potential impact.}
Fraud-risk estimates can generate harmful false accusations. Evaluation
should quantify false positives and calibration, examine subgroup effects
where ethically and statistically appropriate, and treat inferred onset as
uncertain model output. A research score alone should not be interpreted as
evidence sufficient for automatic punitive action.

\paragraph{Working conclusion.}
Separating occurrence from detection gives an explicit account of hidden
fraud among undetected users. Whether it yields reliable onset inference or
better early warning remains an empirical and statistical question. The
immediate goal is to determine what can be recovered under controlled delay
and supervision assumptions before expanding the architecture.

\bibliographystyle{named}
\bibliography{references}
\end{document}
